\subsection{Что чему соответствует}

Во всех трёх системах роли примерно одни и те же: кто-то ведёт учёт ресурсов на узлах, кто-то решает, куда поставить задачу, агент на узле запускает процесс и следит за ним, и кто-то знает логику самого приложения (какие задачи, в каком порядке, что делать при отказе). Отличаются системы тем, кому какая роль достаётся.

\tabin{mapping}

Больше всего на моих прогонах ощущалась разница в последней роли. В YARN у приложения MapReduce есть свой ApplicationMaster. Он знает про фазы map и reduce, сам запрашивает контейнеры у ResourceManager, перезапускает упавшие задачи и запускает reduce, когда закончился map. В Kubernetes такого слоя нет. Job-контроллер держит нужное число успешных подов и перезапускает упавшие индексы, но про то, что reduce можно начинать только после всех map, он не знает. У меня эту работу делает скрипт \mono{scripts/05\_parallel.ps1}: ждёт \mono{par-map}, отправляет \mono{par-reduce}, потом \mono{par-commit}. В настоящих системах для этого ставят Argo Workflows, оператор Spark или Kubeflow, то есть надстройки. В Mesos ещё дальше: сам Mesos задачи не размещает, а только предлагает ресурсы, и размещение, и порядок пишет фреймворк.

\subsection{Кто решает, где запустить задачу}

\textbf{YARN.} Планировщик один, он в ResourceManager. ApplicationMaster присылает запросы вида «$N$ контейнеров по $X$\,МБ и $Y$ vcores, лучше на таком-то узле или стойке». Планировщик (Capacity или Fair) выдаёт контейнеры по мере освобождения ресурсов; о свободных ресурсах он узнаёт из пульса NodeManager. Место выбирает RM, но с учётом пожеланий приложения о локальности данных.

\textbf{Mesos.} Планирование двухуровневое. Мастер не знает, что нужно задачам. Его аллокатор (DRF, Dominant Resource Fairness между ролями) \emph{предлагает} фреймворку ресурсы конкретного агента, а планировщик фреймворка решает, взять их и что на них запустить или отказаться. Выбор места целиком на фреймворке, зато мастер простой.

\textbf{Kubernetes.} Планировщик тоже один (\mono{kube-scheduler}), и о приложении он ничего не знает. Он берёт поды из очереди по одному (очередь отсортирована по приоритету), отбрасывает узлы, куда под не помещается (\mono{requests} против \mono{allocatable}, \mono{nodeSelector}, taints), оценивает оставшиеся и записывает \mono{nodeName}. В моих манифестах это видно напрямую: поды map с \mono{nodeSelector: hw5/pool=batch} попадали только на два воркера, а RustFS с допуском к taint'у control-plane --- только на \mono{hw5-control-plane}. Индексы map планировщик разложил по воркерам примерно поровну (\parNodeSplit\ подов), хотя о сплитах ничего не знает.

\subsection{Как задаются ресурсы и что бывает при превышении}

\textbf{YARN.} У контейнера один размер: сколько памяти и vcores запросило приложение, с округлением вверх до кратного \mono{yarn.scheduler.minimum-allocation-mb} (по умолчанию 1024\,МБ). Отдельно «гарантии» и «потолка» нет. Память проверяет NodeManager: монитор контейнеров раз в несколько секунд складывает RSS всего дерева процессов и, если сумма больше выделенного, убивает контейнер с сообщением \mono{is running beyond physical memory limits}. Есть и проверка виртуальной памяти (\mono{vmem-pmem-ratio}), на которой контейнеры чаще всего и падают. Строгий предел через cgroups тоже можно включить. Процессор в Capacity Scheduler по умолчанию вообще не учитывается: \mono{DefaultResourceCalculator} смотрит только на память, vcores начинают учитываться с \mono{DominantResourceCalculator}, а жёстко ограничиваются только через cgroups.

\textbf{Kubernetes.} Величин две: \mono{requests} учитывает планировщик (сколько «занято» на узле), \mono{limits} kubelet записывает в cgroup контейнера (\mono{memory.max}, \mono{cpu.max}). Память проверяет не kubelet по таймеру, а ядро в момент выделения страницы, поэтому OOM наступает сразу и «проскочить» между замерами нельзя. В моём запуске с пределом 256Mi контейнер прожил около \oomSeconds\,с: ядро сначала выбросило кэш файлов, потом освобождать стало нечего, и оно вызвало OOM-killer внутри cgroup пода. Kubelet ставит \mono{memory.oom.group=1}, поэтому были убиты все процессы контейнера. В YARN при превышении тоже убивается весь контейнер, так что смысл одинаковый, а механизмы разные: опрос RSS по таймеру против проверки в ядре. Процессор Kubernetes ограничивает квотой CFS. У расчёта из ДЗ 3 с \mono{cpu.max} в 4 CPU квота кончалась в \hwThrPeriods\ периодов по 100\,мс: в этих периодах процесс ждал следующего, уже будучи запущенным.

\textbf{Mesos.} Ресурсы приходят в предложениях (\mono{cpus}, \mono{mem}, \mono{disk}, порты). Задача забирает часть предложения, а изоляцию на агенте делает containerizer через те же cgroups.

\subsection{Что происходит, когда ресурсов не хватает}

В \textbf{YARN} запросы контейнеров от ApplicationMaster просто остаются невыполненными: приложение в состоянии \mono{RUNNING}, часть задач ждёт. Если в очереди нет места даже под сам ApplicationMaster, приложение ждёт в состоянии \mono{ACCEPTED}. В \textbf{Mesos} фреймворк не получает подходящих предложений (или отказывается от неподходящих) и ждёт. В \textbf{Kubernetes} под создаётся сразу и висит в \mono{Pending} с событием \mono{FailedScheduling}. В моём прогоне с 1500m CPU на под планировщик написал \mono{Insufficient cpu} для обоих воркеров, хотя на каждом оставалось по 0,9 CPU: два пода по 1,5 CPU занимают 3 из 3,9, третий не помещается. Это обычная фрагментация при упаковке: в YARN она тоже есть и даже крупнее, потому что vcores там целые, а память округляется до \mono{minimum-allocation-mb}. Дольше всех под ждал \cpuPendingMax\,с, а map в целом шёл \cpuMapSpan\,с вместо \parMapSpan.

\subsection{Приоритеты и вытеснение}

В \textbf{YARN} (Capacity Scheduler) вытеснение в первую очередь работает между очередями. Если очередь заняла больше гарантированной доли, а другой очереди её доля нужна, монитор вытеснения снимает контейнеры у той, что взяла лишнее. В Hadoop 3 есть и вытеснение внутри очереди по приоритету приложений (\mono{intra-queue-preemption}).

В \textbf{Mesos} общего вытеснения между фреймворками нет. Мастер может перестать предлагать ресурсы или попросить их вернуть (inverse offers, для обслуживания узлов). Есть revocable-ресурсы: выделенные, но фактически не используемые ресурсы, которые раздаются повторно (переподписка) и которые можно отобрать. Что именно остановить, решает фреймворк. Проект при этом почти не развивается: последний релиз 1.11.0 вышел в 2020 году.

В \textbf{Kubernetes} вытеснение идёт по приоритету пода. Если под с высоким \mono{PriorityClass} никуда не помещается, планировщик ищет узел, где удаление подов с меньшим приоритетом освободит место, записывает поду \mono{nominatedNodeName} и удаляет выбранные поды. В моём прогоне срочный расчёт (\mono{batch-high}, 2 CPU) пришёл, когда поды \mono{batch-low} занимали на воркерах все места. Планировщик снял \preemptedCount\ пода на одном узле, и они получили условие \mono{DisruptionTarget} с причиной \mono{PreemptionByScheduler}. Чтобы фоновое задание из-за этого не упало, в манифесте map есть правило \mono{podFailurePolicy} с действием \mono{Ignore} для \mono{DisruptionTarget}. Вытесненные индексы создаются заново и не тратят \mono{backoffLimitPerIndex}. Фоновый расчёт закончился позже, но с тем же результатом, что в ДЗ 3. Повторно считать вытесненный сплит безопасно: map пишет вывод сплита в S3 целиком и только потом ставит \mono{\_SUCCESS}.

\subsection{Квоты}

Очередь Capacity Scheduler в \textbf{YARN} --- это доля кластера с гарантией (\mono{capacity}) и потолком (\mono{maximum-capacity}). Пока другие очереди простаивают, очередь может занять больше гарантии, а когда ресурсы нужны другим, её поджимают вытеснением. В \textbf{Mesos} похожую роль играют роли с весами (доля в DRF), квотами и резервированиями ресурсов на агентах.

\textbf{ResourceQuota} в Kubernetes устроена иначе. Это жёсткий потолок на сумму \mono{requests} и \mono{limits} в пространстве имён, и проверяется он при \emph{создании} пода, а не при планировании. Поэтому в \mono{hw5-quota} лишние поды map не висели в \mono{Pending}: их просто не удавалось создать. Job-контроллер получал \mono{FailedCreate: exceeded quota} и пробовал снова при следующей синхронизации, и поды шли волнами по три, по мере того как заканчивались предыдущие, хотя на воркерах оставалось свободно ещё \quotaFreeCpu\ CPU: занять их квота не даёт. Гарантии, как в YARN, она тоже не даёт: если чужие поды займут кластер, мои будут ждать при неиспользованной квоте. Очереди с заимствованием и справедливым делением в Kubernetes добавляют надстройки: Kueue (\mono{ClusterQueue} и cohort), Volcano, YuniKorn (последний прямо переносит модель очередей YARN). Ещё одно отличие: квота на \mono{cpu} и \mono{memory} требует, чтобы эти поля были у каждого пода. Под без них отклоняется, пока нет \mono{LimitRange} со значениями по умолчанию; в YARN такой проблемы нет, там всегда есть минимальный размер контейнера.

\subsection{Что когда выбирать}

YARN рассчитан на пакетную аналитику вокруг HDFS: локальность данных, очереди с долями между командами, готовые ApplicationMaster для MapReduce, Spark и Tez. Mesos --- тонкий слой распределения ресурсов, поверх которого каждый фреймворк пишет свой планировщик, гибко, но всё приходится делать самому. Kubernetes держит описанное состояние подов и одинаково подходит и для сервисов, и для заданий, а пакетные вещи (стадии, очереди с заимствованием, gang scheduling) приходят надстройками. В моей задаче данные лежат в S3, а не в HDFS, так что главное преимущество YARN, локальность, здесь не работает. Kubernetes дал всё, что требовалось: пределы памяти с OOM в ядре, параллельное задание по индексам, приоритеты и квоты. Связку стадий map $\to$ reduce пришлось делать скриптом.
